\begin{frame}{Reuse saves resources when avoided preparation exceeds its overhead.\ev{\tagA}}
\vspace{.3cm}
\begin{center}
{\fontsize{22}{26}\selectfont$(N-1)P > H + N(R+D)$\par}
\vspace{.55cm}
{\fontsize{12.5}{15}\selectfont
$P$: preparation per continuation\qquad $H$: capture\par
\vspace{.18cm}
$R$: restore\qquad $D$: incremental divergence\qquad $N$: children\par
\vspace{.35cm}
Additive resource cost; common continuation work cancels.\par}
\end{center}
\sources{Illustrative resource-cost model; Jitsu, NSDI 2015; LightVM, SOSP 2017.}
\qadetail{This is an additive resource-cost model, not measured elapsed time. The ordinary continuation work common to both alternatives cancels. D is incremental divergence overhead, including copy-on-write. P must be the reusable preparation rather than merely sandbox setup. Resource cost and batch elapsed time are different endpoints: elapsed time depends on critical paths and contention. The original numerical cases are retained in the next slide: N8, P30, H2, R1, D3 avoids210 resource-seconds and adds34, saving176; replacing P30 with a warm-cache P2 avoids14 and adds34, losing20. In SELFHOST-2 the three serial edit tasks sum to342.579 seconds, versus an ideal unchanged-duration critical path176.689 seconds, a computed165.890-second saving or about6.5 percent of recorded stage time, before fork, merge, contention and verification overhead. Nothing forked in that work order; worktrees may support the file-level schedule. Jitsu and LightVM illustrate cheap fresh-start alternatives.}
\note{[0:45] The resource comparison has a simple boundary. Creating children avoids repeating preparation, but pays capture, restore and incremental divergence costs. Work that both alternatives perform cancels. The remaining question is whether avoided preparation exceeds that overhead. Preparation means the state a continuation can actually reuse. Our twenty-second sandbox setup is only one part of a possible prefix. This model counts resources; batch elapsed time also depends on scheduling and contention. The next calculation shows why a good warm alternative can change the answer.}
\end{frame}
